\section{Scope and limitations}\label{sec:scope}

The result is an exact infinite-data information theorem.  It assumes a
site-pattern tensor generated exactly by the stated model and classifies the
semi-directed topology, not all numerical parameters.  In particular,
incidence gauges in \cref{thm:bridge-fibre} prevent the topology theorem from
identifying individual physical bridge multipliers.

The theorem is restricted to:

\begin{itemize}
  \item binary standard semi-directed networks under the fixed one-step
        root-suppression convention;
  \item level at most two;
  \item strong tree-childness for the classification;
  \item positive K3P Fourier eigenvalues and strictly positive transition
        probabilities;
  \item inheritance probabilities strictly between zero and one;
  \item fixed observable state labels \(C,G,T\); and
  \item the principal positive domain, or separately the strict
        continuous-time cone.
\end{itemize}

It does not cover nonbinary networks, higher level, arbitrary weakly
tree-child networks, boundary probabilities, zero or signed eigenvalues, an
untransported permutation of the three nonzero states, or arbitrary
nonreversible substitution models.  No claim is made for the signed K3P
components outside \(\Dplus\).

Generic identifiability means outside a proper algebraic exceptional set for
each fixed topology.  It is not pointwise topology identifiability at every
parameter, equality of complete stochastic images, or universal
identifiability of physical parameters.  The exact reconstruction theorem is
a termination result under an exact-real oracle.  It supplies no finite-sample
estimator, statistical error rate, numerical conditioning bound, or practical
sequence-length guarantee.

The weak-class construction proves sharpness of the hypothesis in the sense
that weakening ``strongly'' to ``weakly'' makes the classification false.  It
does not classify every weakly tree-child network or every possible weak-class
ambiguity.
